\documentclass[10pt]{article}
\providecommand{\CRevisionRound}{2}
\usepackage[margin=0.78in]{geometry}
\usepackage{fontspec}
\setmainfont{TeX Gyre Pagella}
\setsansfont{TeX Gyre Heros}
\setmonofont{Latin Modern Mono}
\usepackage{microtype,amsmath,amssymb,amsthm,mathtools,booktabs,array,enumitem}
\usepackage[hidelinks]{hyperref}
\pdfvariable suppressoptionalinfo 767
\setlength{\parindent}{0pt}
\setlength{\parskip}{4pt}
\newtheorem{theorem}{Theorem}
\newtheorem{lemma}{Lemma}
\title{Cut and First-Conjugate Geometry of the Standard Heisenberg Group:\\A Complete Hamiltonian and Distance Atlas}
\author{HCS Research Program}
\date{1 September 2026}

\begin{document}
\maketitle

\begin{abstract}
For the standard real Heisenberg group we integrate every unit-speed normal
sub-Riemannian geodesic and close the global minimizing geometry, rather than
stopping at an explicit orbit formula.  We exclude nonconstant abnormal
extremals, compute the exponential Jacobian, prove equality of the first cut,
conjugate, and rotational Maxwell times, and determine the cut locus from the
identity.  A strictly monotone implicit angle gives the
Carnot--Carath\'eodory distance, including both degenerate endpoint faces.  All
signs and factors are frozen to one frame convention; no extension to
arbitrary Carnot groups is claimed.
\end{abstract}

\section{Frozen model}
On $H^1=\mathbb R^3$ put
\[
 X=\partial_x-\frac y2\partial_z,\qquad
 Y=\partial_y+\frac x2\partial_z,\qquad [X,Y]=\partial_z,
\]
and make $X,Y$ orthonormal.  The normal Hamiltonian is
\[
 H=\frac12\left(p_x-\frac y2p_z\right)^2+
   \frac12\left(p_y+\frac x2p_z\right)^2.
\]
Write $\lambda=p_z$ and impose $H=1/2$.

\begin{theorem}[Complete identity-based atlas]
Every unit-speed normal geodesic from the identity has an angle $\phi$ and is,
when $\lambda=0$,
\[
 (x,y,z)=(t\cos\phi,t\sin\phi,0),
\]
while for $\lambda\ne0$ it is
\begin{align*}
 x&=\frac{\sin(\phi+\lambda t)-\sin\phi}{\lambda},&
 y&=\frac{\cos\phi-\cos(\phi+\lambda t)}{\lambda},\\
 z&=\frac{\lambda t-\sin(\lambda t)}{2\lambda^2}.
\end{align*}
There are no nonconstant abnormal extremals.  The first conjugate and cut
times agree:
\[
 t_{\rm conj}=t_{\rm cut}=\frac{2\pi}{|\lambda|}\quad(\lambda\ne0),
\]
and the cut locus and first conjugate locus from the identity are both
$\{(0,0,z):z\ne0\}$.

For $\rho=(x^2+y^2)^{1/2}>0$, let $\theta\in(-\pi,\pi)$ be the unique solution
of
\[
 \frac{4z}{\rho^2}=\frac{\theta-\sin\theta\cos\theta}{\sin^2\theta}.
\]
Then
\[
 d(0,(x,y,z))=\rho\frac{\theta}{\sin\theta}.
\]
Continuously, $d=\rho$ when $z=0$, whereas
$d(0,(0,0,z))=2\sqrt{\pi|z|}$ on the vertical axis.
\end{theorem}

\section{Hamilton integration and abnormal boundary}
Set $h_1=p(X)$ and $h_2=p(Y)$.  Hamilton's equations give
\[
 \dot h_1=-\lambda h_2,\qquad \dot h_2=\lambda h_1,
 \qquad \dot\lambda=0.
\]
Thus $(h_1,h_2)=(\cos(\phi+\lambda t),\sin(\phi+\lambda t))$; integrating
$\dot x=h_1$, $\dot y=h_2$ and
$\dot z=(-yh_1+xh_2)/2$ proves the displayed formulas.  An abnormal covector
annihilates $X,Y$ and, by differentiating these constraints, $[X,Y]$.
Since these three vectors span $TH^1$, the covector vanishes, a contradiction.

\section{First singular and Maxwell times}
For initial horizontal norm $r$, direct differentiation in
$(r,\phi,\lambda)$ gives, with $s=\lambda t$,
\[
 \det D\operatorname{Exp}_t=
 \frac{r^3t}{\lambda^4}\{2-2\cos s-s\sin s\}.
\]
The brace factors as
$4\sin(s/2)[\sin(s/2)-(s/2)\cos(s/2)]$.  For $0<s<2\pi$ the first factor is
positive; if $u=s/2$, the derivative of $\sin u-u\cos u$ is $u\sin u>0$.
Hence the first zero is $s=2\pi$.  At that phase $x=y=0$ and
$z=\operatorname{sgn}(\lambda)\pi/\lambda^2$, independently of $\phi$.
The conjugate time is therefore also the first rotational Maxwell time.

\ifnum\CRevisionRound>0
\begin{lemma}[Convention audit at zero and negative momentum]
The determinant formula extends continuously to $\lambda=0$ with value
$r^3t^5/12$.  For $\lambda<0$ its first zero is still
$2\pi/|\lambda|$, and the corresponding endpoint has negative $z$.
\end{lemma}
\begin{proof}
The expansion $2-2\cos s-s\sin s=s^4/12+O(s^6)$ gives the first assertion.
The brace is even in $s$, whereas
$(s-\sin s)/(2\lambda^2)$ is odd in $s$; this proves both remaining sign
statements.  Thus neither a hidden singular zero-momentum face nor a sign
choice halves the conjugate time.
\end{proof}
\fi

\section{Distance and minimality}
The length of a horizontal curve equals the Euclidean length of its
$(x,y)$ projection, while
\[
 z=\frac12\int(x\,dy-y\,dx)
\]
is signed planar area.  Dido's problem says that a minimizing projected curve
with fixed endpoint and signed area is a line or circular arc.  Before a full
turn the sub-full-turn Dido arc is unique; at a full turn its starting tangent
is free.
Writing its half-angle as $\theta$ yields
\[
 \rho=\frac{2\sin\theta}{\lambda},\qquad
 z=\frac{\theta-\sin\theta\cos\theta}{\lambda^2},\qquad
 t=\frac{2\theta}{\lambda},
\]
and hence the theorem.  Moreover
\[
 \frac{d}{d\theta}\frac{\theta-\sin\theta\cos\theta}{\sin^2\theta}
 =\frac{2(\sin\theta-\theta\cos\theta)}{\sin^3\theta}>0
\]
on $(-\pi,\pi)$, proving angle uniqueness.  The isoperimetric inequality for a
closed planar curve gives $t^2\ge4\pi|z|$, with equality for a circle, proving
the vertical formula and completing the cut-locus argument.

\ifnum\CRevisionRound>1
\begin{lemma}[Why the Maxwell set is the cut locus]
No geodesic segment with $|\lambda|t<2\pi$ loses global minimality, and every
segment with $|\lambda|t>2\pi$ is nonminimizing.
\end{lemma}
\begin{proof}
The planar constrained-length Euler equation has constant signed curvature,
so every nonstraight minimizer is a circular arc.  Strict monotonicity of the
implicit-angle map selects exactly one arc with half-angle in $(-\pi,\pi)$ for
each nonvertical endpoint.  Hence no competitor exists before the limiting
full circle.  At the full circle the isoperimetric equality gives a minimizing
$S^1$ family.  Past it, the same endpoint constraints select the unique
sub-full-turn Dido arc whose half-angle lies in $(-\pi,\pi)$, so the continued
geodesic is longer.
This upgrades the visible Maxwell merger to the claimed global cut theorem.
\end{proof}

The three endpoint faces fit without an extra normalization:
\[
\begin{array}{c@{\qquad}c@{\qquad}c}
z=0,\ \rho>0 & \rho>0,\ z\ne0 & \rho=0,\ z\ne0\\
\theta=0 & -\pi<\theta<\pi & \theta\to\operatorname{sgn}(z)\pi\\
d=\rho & d=\rho\theta/\sin\theta & d=2\sqrt{\pi|z|}.
\end{array}
\]
\fi

\section{Evidence and claim boundary}
The deterministic receipt contains 800 trajectory rows, 64 regular-distance
rows, and 12 vertical rows (10,972 numeric cells, independently recounted from
an explicit field schema).  An implementation sharing no producer code passes
11,139 assertions; 20 symbolic identities, byte replay, and 27/27 repaired-hash
hostile mutations also pass.  These are regression
checks: the global theorem is established by the analytic argument above.

\ifnum\CRevisionRound>1
\begin{center}
\small
\begin{tabular}{@{}lr@{}}
\toprule
Receipt or gate & Result\\
\midrule
Trajectory / distance / vertical rows & $800/64/12$\\
Independent assertions & $11{,}139$ PASS\\
Symbolic identities & $20$ PASS\\
Repaired-hash hostile mutations & $27/27$ rejected\\
Evidence SHA-256 & \texttt{86bbae2b0610\ldots dcaa}\\
\bottomrule
\end{tabular}
\end{center}
The producer replay is byte-identical.  The independent checker imports no
producer code, while the symbolic checker derives the decisive identities from
the frozen frame.  Infinite-dimensional or global conclusions are therefore
not inferred from a finite numerical grid.
\fi

Complete geodesics provide no periodic-orbit signal: $\lambda=0$ gives lines,
while for $\lambda\ne0$ one horizontal period has nonzero vertical drift.
Thus the A1 axis is conservatively failed.  The theorem status is
\textbf{PROVABLE AS STATED}.  The Route-A tuple is
\begin{center}\small
\begin{tabular}{@{}c@{}}
\texttt{(A0\_FAIL, A1\_FAIL, A2\_FAIL, A3\_FAIL,}\\
\texttt{A4\_FORMAL\_HINT)}, hence \texttt{ROUTE\_A\_REJECTED}.
\end{tabular}
\end{center}
Route B is not invoked.  The claim firewall is
\texttt{NO\_BAD\_EULER\_OR\_ROOT\_NUMBER}: no arithmetic local data, Euler
factor, root number, automorphy, target divisor, or Hilbert--P\'olya operator is
claimed.

\section*{Primary source}
B. Gaveau, ``Principe de moindre action, propagation de la chaleur et estim\'ees
sous elliptiques sur certains groupes nilpotents,'' \emph{Acta Mathematica}
\textbf{139} (1977), 95--153. DOI:
\href{https://doi.org/10.1007/BF02392235}{10.1007/BF02392235}.
The citation records classical model lineage; no priority claim is made.

\end{document}
